% ---
% id: meat-core/thm:keylength
% kind: thm
% title: "Secrecy of the key from the accumulated entropy"
% summary: >-
%   For a Rényi order α in (1, 2), a key of length ℓ ≤ n·h_low − λ_EC −
%   ⌈log₂(1/ε_EV)⌉ − α/(α−1)·log₂(1/ε_PA) + 1 is ε_PA-secret on the accepted
%   state at every all-at-once finite-register attack the reduction chain
%   carries onto the round model ([[A-LINK]]). [[cor:composed-security]]
%   extends it to every attack.
% user_statement: >-
%   Over n rounds the attacker's uncertainty about the raw key accumulates
%   to n times the per-round bound h_low, less a term for how unlikely the
%   acceptance test is to pass, which cancels over all runs; the bound h_low
%   holds over every attack compatible with your declared devices and with
%   passing statistics. Subtracted are the
%   error-correction transcript of λ_EC bits, the verification hash of
%   ⌈log₂(1/ε_EV)⌉ bits and a finite-size penalty of α/(α−1)·log₂(1/ε_PA)
%   bits; one bit is added back. The remaining key of ℓ bits is secret,
%   counted over all runs: it is within the secrecy parameter ε_PA of a
%   uniform key independent of the attacker's memory and every public
%   message. It is proved for a finite-memory attacker receiving every
%   signal at once; [[cor:composed-security]] covers every attacker.
%   The Rényi order α in (1, 2) is fixed before the run ([[A-FIXED]]) and
%   changes the key's length, never its secrecy.
% prominence: main
% status: implemented
% uses: [lem:memory-lift, def:postprocessing-model, thm:accumulated-entropy, def:homog, thm:reduction, cor:reduction-constrained, def:feasible, def:roundmodel, thm:constrained-accumulation, foundations/def:renyi, def:admissible-declaration]
% mentions: [cor:constrained-key-length, rem:plus2, def:authenticated-channel, cor:composed-security, protocol-compiler/thm:band-a, protocol-compiler/rem:obligations-discharged, reductions/thm:two-sided, reductions/cor:sequential, reductions/thm:licence]
% assumes: [A-KEYREG, A-HASH, A-ATK, A-SEED, A-FIXED, A-CONFORM, A-LAB, A-SEQ, A-LINK, A-ENS, A-SPEC, A-ALICE-MEAS]
% refs:
%   - "tomamichel2016finite, Lemmas 5.3-5.4, Eq. (5.94)"
%   - "tomamichel2016finite, Thm. 4.1, Eq. (4.82)"
%   - "kamin2025renyi, Thm. 4, Eq. (13)"
%   - "tupkary2026decoy, Lemma 6.5, Rem. 6.1"
% provenance: newly-derived
% implements: certiqs_key.meat_core.finitekey.keylength:key_length
% lean:
% notes: "The acceptance flag of error verification is part of the public transcript, so secrecy holds on the accepted state at the cost of one bit: the constant +2 of the derivation this theorem generalises becomes +1 here, and with that derivation's weaker convention (secrecy conditioned on the acceptance test only) its constant is recovered (rem:plus2). The side information E_tot holds the public seeds S; they cost nothing, being independent of the raw key and the adversary's other registers."
% ---

The accumulated entropy bound feeds the fixed-length post-processing model:
error correction with a transcript of fixed length, error verification whose
one-bit outcome is announced, and privacy amplification by an ideal
two-universal hash. Announcing the verification outcome costs one bit and
makes the secrecy statement hold on the accepted state
(Remark~\ref{rem:plus2}).

The theorem is stated at the attacks the analysis is carried out for: those
with finite registers whose first channel receives every transmitted signal,
which the \tref{reduction chain} carries onto the \tref{round model}, as
hypothesis (H-CHAIN) says. Its link part is proved, not attested, on the
identity link of \tref{Band A} (\xref{protocol-compiler/thm:band-a}) and for
a chain the reductions module emits with status exact
(\xref{reductions/thm:two-sided}, \xref{reductions/cor:sequential},
\xref{reductions/thm:licence}), where A-LINK is discharged
(\xref{protocol-compiler/rem:obligations-discharged}); for any other chain
it rests on A-LINK. With constraints, the part that makes the attacks
admissible for the constraint list rests on A-ATK on every route, Band A
included, wherever the list carries an entry no theorem licenses.
Corollary~\ref{cor:composed-security} extends the bound to every attack,
whatever the adversary's memory.

\begin{theorem}\label{thm:keylength}
Let the protocol run the \tref{fixed-length post-processing model}
(Definition~\ref{def:postprocessing-model}) in the
\tref{honest-authentication setting}, with the registers
\eqref{eq:pp-registers}. Let one of the two accumulation branches hold:
\emph{(unconstrained)} the hypotheses of
Theorem~\ref{thm:accumulated-entropy} with $\CList=\varnothing$; or
\emph{(constrained)} the hypotheses of
Theorem~\ref{thm:constrained-accumulation} ($\CList$
\tref{licensed for the declared source mode}, Definition~\ref{def:homog};
A-ATK (\xref{protocol-compiler/ass:atk})), in either branch save the
condition $\Pr[\OmAT]>0$ that the accumulation theorem places on the
attack, which the attack below need not meet. Let the run be as in
Lemma~\ref{lem:memory-lift}\,(a), as A-SPEC (\xref{protocol-compiler/ass:spec})\,(i)
attests for its measurement, with A-ALICE-MEAS
(\xref{protocol-compiler/ass:alice-meas})\,(i) for Alice's measurement where
no \tref{source-replacement certificate} is demanded, and A-ENS
(\xref{transmitter/ass:ens}), or the
premise a \tref{source model} declares in its place, for its source states;
let it be attacked by an attack in the sense of
Lemma~\ref{lem:memory-lift}\,(b) that is all-at-once and has finite
registers; and let the \tref{reduction chain} from which the declaration was
compiled carry it onto the round model, as A-LINK
(\xref{protocol-compiler/ass:link})\,(i)(a), (i)(b), (ii)(a) and (iv) attest,
the existence of each link's channel, the pull-back identity, the source
identity and complete positivity, with A-ATK in the constrained branch:
\begin{itemize}[leftmargin=2.4em]
\item[(H-CHAIN)] the marginal on $Z_1^n\Chat_1^nI_1^nE_n$ of the run's state
  after the $n$ rounds is the state the
  \tref[round channel]{round channels} of Definition~\ref{def:roundmodel}
  produce under an \tref{attack sequence} of that definition,
  $\CList$-admissible in the constrained branch, from an input state as in
  the branch's accumulation theorem; the run's state is an extension of it,
  holding the parties' further records.
\end{itemize}
Let Assumptions~\ref{ass:keyreg}, \ref{ass:hash}, \ref{ass:seed},
\ref{ass:fixed}, \ref{ass:conform} and~\ref{ass:lab}, the last at each of
the two parties, and A-SEQ (\xref{transmitter/ass:seq}) hold, let
$\alpha\in(1,2)$ and $\epa,\eev\in(0,1]$. If the output length satisfies
\begin{equation}\label{eq:keylength}
  \ell\;\le\;n\,\hlow-\lEC-\Big\lceil\log_2\frac1{\eev}\Big\rceil
  -\frac{\alpha}{\alpha-1}\log_2\frac1{\epa}+1,
\end{equation}
then, under that attack, the protocol is $\epa$-secret on the accepted
state:
\begin{equation*}
  \frac12\Big\|\rho_{K\,E_{\mathrm{tot}}\wedge\Omega_{\mathrm{acc}}}
  -\frac{\idon{K}}{2^{\ell}}\otimes
  \rho_{E_{\mathrm{tot}}\wedge\Omega_{\mathrm{acc}}}\Big\|_1\le\epa .
\end{equation*}
\end{theorem}

\begin{proof}
\emph{Idea.} Secrecy is first proved on the state conditioned on the
acceptance test: the chain rule for a classical register charges the whole
public transcript $L$ against the accumulated entropy, the public seeds cost
nothing because they are independent of everything else, and the R\'enyi
leftover-hashing lemma turns the result into a trace-distance bound, whose
acceptance-probability factor cancels. A contraction then passes from the
acceptance test to the accepted state, on which the announced flag
$F_{\mathrm{acc}}$ is a register of the adversary's side information. An
attack under which the test never passes is disposed of first.

\emph{Step 0 (an attack that never passes the test).} The event $\OmAT$ is
a function of $\Chat_1^n$, so by (H-CHAIN) its probability under the
attack is the one in the round model. If $\Pr[\OmAT]=0$, then
$\Omega_{\mathrm{acc}}=\OmAT\wedge\{F_{\mathrm{acc}}=\mathrm{acc}\}$
(Definition~\ref{def:postprocessing-model}) has probability zero, so
$\rho_{K\,E_{\mathrm{tot}}\wedge\Omega_{\mathrm{acc}}}=0$ and
$\rho_{E_{\mathrm{tot}}\wedge\Omega_{\mathrm{acc}}}=0$, and the left-hand
side of the secrecy criterion is $0\le\epa$. From now on $\Pr[\OmAT]>0$,
and the hypotheses of the branch's accumulation theorem hold in full.

\emph{Step 1 (secrecy on the acceptance test).} Abbreviate
$B\coloneqq\Chat_1^nI_1^n\Ehat$. By (H-CHAIN) the rounds of the run under the
attack leave, on $Z_1^n\Chat_1^nI_1^nE_n$, the state the
\tref[round channel]{round channels} of Definition~\ref{def:roundmodel}
produce under an attack sequence of that definition, to which the
accumulation theorem of the branch applies, the run's further records being
an extension of it, and
the register $E_{\mathrm{tot}}=B\,S\,L\,G$ is what the adversary holds at the
end of the protocol. By A-SEQ (\xref{transmitter/ass:seq})\,(i), (ii)
and~(iv), every message of a round is sent after every measurement of that
round, no operation on a quantum system of a round depends on that round's
messages, and the classical post-processing is applied after the $n$
rounds, so each round is the measurement of its round channel followed by
announcements that the model hands to the adversary, and the post-processing
acts on the state the $n$ round channels produce. By
Assumption~\ref{ass:lab} at each of the two parties, during the protocol
nothing about the raw string, the keys or the parties' local randomness
leaves either laboratory, other than as the protocol's output, except
through the systems the \tref{round model} hands to the adversary and the
public messages, and the adversary acts on neither laboratory except through
the systems and messages the model routes to it. By
Assumption~\ref{ass:conform}\,(ii) the classical messages after the rounds
are exactly those of the model and of the authentication post-processing
(Definition~\ref{def:authenticated-channel}\,(c)). In the
honest-authentication setting no message fails authentication, so the
messages of the authentication post-processing are accept or abort
decisions that are functions of $\Omega_{\mathrm{acc}}$, hence of
$\Chat_1^n$ and $F_{\mathrm{acc}}$, which $E_{\mathrm{tot}}$ holds;
appending them is a channel acting on $E_{\mathrm{tot}}$ alone, which does
not increase the trace distance of the secrecy criterion and keeps its ideal
term of the same form, so they cost nothing. By
Assumption~\ref{ass:fixed}, $n$, $\acSet$, $\ell$, $\lEC$, $\epa$ and $\eev$
are constants of the run, so the event $\OmAT$, the dimension $2^{\lambda}$
of $L$ and the output length do not depend on the run's data, and by
Assumption~\ref{ass:conform}\,(i) and~(iv) the run applies that test and
outputs the $\ell$-bit key the model computes, or $\perp$. The
post-processing is executable: the key holder computes $Z_1^n$ from its own
readout and the announcements (Definition~\ref{def:admissible-declaration}\,(viii)), and every
object the post-processing forms is a function of registers present in the
state that Theorem~\ref{thm:accumulated-entropy} bounds, together with the
public seeds.

The state after the public phase, conditioned on $\OmAT$, is classical on
$L$; by the chain rule for conditioning on a classical register
\cite[Lemmas~5.3--5.4, Eq.~(5.94)]{tomamichel2016finite}, applied with
the classical system $L$ of dimension $2^{\lambda}$ and the conditioning
system $B\,S$ (the $\alpha$-domain
restriction to $(1,2)$ is this module's, inherited from the
leftover-hashing lemma applied next; the cited chain rule is stated for the
sandwiched family of \cite[Def.~5.2, Eqs.~(5.18)--(5.19)]%
{tomamichel2016finite} on its data-processing range
$\alpha\in[\frac12,\infty]$, the range on which
\cite[Thm.~4.1, Eq.~(4.82)]{tomamichel2016finite} grounds data
processing, and $(1,2)\subset[\frac12,\infty]$),
\begin{equation}\label{eq:chainrule}
\begin{aligned}
  \Hsandup\big(Z_1^n\big|B\,S\,L\big)_{\rho_{|\OmAT}}
  &\;\ge\;\Hsandup\big(Z_1^n\big|B\,S\big)_{\rho_{|\OmAT}}-\lambda
  \;\ge\;\Hsandup\big(Z_1^n\big|B\big)_{\rho_{|\OmAT}}-\lambda\\
  &\;\ge\;n\,\hlow-\frac{\alpha}{\alpha-1}\log_2\frac{1}{\Pr[\OmAT]}
  -\lambda .
\end{aligned}
\end{equation}
The second inequality holds because the seeds cost nothing. By
Assumption~\ref{ass:seed} they are drawn after the rounds, independently of
every other register, so that
$\rho_{Z_1^nBS|\OmAT}=\rho_{Z_1^nB|\OmAT}\otimes\omega_S$ for the law
$\omega_S$ of $S$, $\OmAT$ being a function of $\Chat_1^n$. For every
$\sigma_B\in\Seq{B}$ the state $\sigma_B\otimes\omega_S$ is admissible in the
infimum that defines $\Hsandup(Z_1^n|BS)$ (\xref{foundations/def:renyi}), and
with $\gamma\coloneqq\frac{1-\alpha}{2\alpha}$, powers taken on supports,
\[
  (\idon{Z}\otimes\sigma_B\otimes\omega_S)^{\gamma}
  (\rho_{Z_1^nB}\otimes\omega_S)
  (\idon{Z}\otimes\sigma_B\otimes\omega_S)^{\gamma}
  =\big((\idon{Z}\otimes\sigma_B)^{\gamma}\rho_{Z_1^nB}
  (\idon{Z}\otimes\sigma_B)^{\gamma}\big)\otimes\omega_S^{1/\alpha};
\]
the trace of the $\alpha$-th power of a tensor product is the product of the
traces, $\tr[(\omega_S^{1/\alpha})^{\alpha}]=\tr\omega_S=1$, and the support
condition holds for the product exactly when it holds for the first factor.
Hence
$\Dsand(\rho_{Z_1^nB}\otimes\omega_S\|\idon{Z}\otimes\sigma_B\otimes\omega_S)
=\Dsand(\rho_{Z_1^nB}\|\idon{Z}\otimes\sigma_B)$, and the infimum over
$\sigma_B$ gives $\Hsandup(Z_1^n|BS)\ge\Hsandup(Z_1^n|B)$.

The third inequality of \eqref{eq:chainrule} uses
Theorem~\ref{thm:accumulated-entropy} with
Theorem~\ref{thm:reduction} ($\hqkd\ge\hlow$) in the unconstrained
branch, and Theorem~\ref{thm:constrained-accumulation} with
Corollary~\ref{cor:reduction-constrained} --- that is,
Corollary~\ref{cor:constrained-key-length} --- in the constrained branch;
here and in the rest of the proof $\hlow$ is \eqref{eq:hlow-def} over the
declared $\Feas$, so in the constrained branch it is the constrained value
$\hlow(\CList)$, and the remainder of the proof consumes only the
displayed bound, supplied by either branch. The chain rule holds for an
arbitrary jointly distributed classical $L$ of dimension $2^{\lambda}$ ---
no functional dependence of $L$ on $Z_1^n$ is required --- so it covers
$F_{\mathrm{acc}}$, which depends also on the partner's reconstruction
$\hat Z_1^n$. The third inequality of \eqref{eq:chainrule} applies to the
post-processed state because
$\Hsandup(Z_1^n|B)$ is a functional of the $(Z_1^n,B)$ marginal only, and
the post-processing leaves that marginal unchanged: the registers traced out
by the \tref[round channel]{round channels} (the partner's outcome and
announcement copies;
Definition~\ref{def:roundmodel}) still exist in the physical protocol as an
extension of the modelled state, and every appended object ---
$S$ drawn independently, $L_{\mathrm{EC}}$ and $T_{\mathrm{EV}}$ from
classical copies of $Z_1^n$ and the public seeds, $F_{\mathrm{acc}}$ from
$T_{\mathrm{EV}}$ and the partner's data (Assumption~\ref{ass:keyreg}) --- is
generated by channels acting trivially on $Z_1^nB$. Hence the accumulation
bound of the branch in force, Theorem~\ref{thm:accumulated-entropy} or
Theorem~\ref{thm:constrained-accumulation}, applies verbatim. Now apply the
R\'enyi leftover-hashing lemma \cite[Thm.~4, Eq.~(13)]{kamin2025renyi}
(stated there for $\alpha\in(1,2)$; this is what confines the theorem to
that domain), under the ideal-family hypothesis of
Assumption~\ref{ass:hash}\,(i) (\cite[Lemma~6.5, Rem.~6.1]{tupkary2026decoy};
a merely two-universal family costs one further bit there) and with $G$
drawn with the law of its family independently of $Z_1^nBSL$
(Assumption~\ref{ass:seed}), to the normalized state $\rho_{|\OmAT}$ with
side information $E\coloneqq B\,S\,L$ and hash register $G$:
\begin{equation*}
  \frac12\Big\|\rho_{K\,EG|\OmAT}-\frac{\idon{K}}{2^{\ell}}\otimes
  \rho_{EG|\OmAT}\Big\|_1
  \;\le\;2^{\frac{2(1-\alpha)}{\alpha}}\,
  2^{\frac{1-\alpha}{\alpha}\left(
  \Hsandup(Z_1^n|BSL)_{\rho_{|\OmAT}}-\ell\right)} .
\end{equation*}
Multiply both sides by $\Pr[\OmAT]$ to pass to the subnormalized
$\wedge\OmAT$ state, and insert \eqref{eq:chainrule}: with
$\kappa\coloneqq\frac{\alpha-1}{\alpha}>0$, the factor
$2^{\kappa\log_2(1/\Pr[\OmAT])}=1/\Pr[\OmAT]$ cancels the prefactor
$\Pr[\OmAT]$, leaving
\begin{equation}\label{eq:secrecy-at}
  \frac12\Big\|\rho_{K\,E_{\mathrm{tot}}\wedge\OmAT}
  -\frac{\idon{K}}{2^{\ell}}\otimes
  \rho_{E_{\mathrm{tot}}\wedge\OmAT}\Big\|_1
  \;\le\;2^{-2\kappa}\,2^{-\kappa\left(n\hlow-\lambda-\ell\right)} .
\end{equation}
Requiring the right side to be at most $\epa$ is equivalent to
$n\hlow-\lambda-\ell\ge\frac1{\kappa}\log_2\frac1{\epa}-2$, i.e.\ to
\eqref{eq:keylength} after inserting
$\lambda=\lEC+\lceil\log_2\frac1{\eev}\rceil+1$.

\emph{Step 2 (from $\OmAT$ to $\Omega_{\mathrm{acc}}$).} The flag
$F_{\mathrm{acc}}$ is a register of $E_{\mathrm{tot}}$ (it is part of
$L$). Let $\Pi\coloneqq\dyad{\mathrm{acc}}_{F_{\mathrm{acc}}}$, acting on
$E_{\mathrm{tot}}$ only. The selection map
$X\mapsto(\idon{K}\otimes\Pi)X(\idon{K}\otimes\Pi)$ is trace-norm
contractive ($\|AXB\|_1\le\|A\|_{\infty}\|X\|_1\|B\|_{\infty}$ with
$\|\Pi\|_{\infty}=1$), maps
$\rho_{\wedge\OmAT}\mapsto\rho_{\wedge\Omega_{\mathrm{acc}}}$, and maps
the ideal term to
$\frac{\idon{K}}{2^{\ell}}\otimes
\Pi\rho_{E_{\mathrm{tot}}\wedge\OmAT}\Pi
=\frac{\idon{K}}{2^{\ell}}\otimes
\rho_{E_{\mathrm{tot}}\wedge\Omega_{\mathrm{acc}}}$, preserving the ideal
form because $\Pi$ does not act on $K$. Applying it to
\eqref{eq:secrecy-at} gives the claimed accepted-state secrecy.
\end{proof}
